\documentclass[11pt]{article}
\usepackage[utf8]{inputenc}
\usepackage{amsmath,amssymb,amsthm}
\usepackage[margin=1in]{geometry}
\usepackage{hyperref}
\usepackage{xcolor}
\usepackage{enumitem}
\usepackage{newunicodechar}
\newunicodechar{ℝ}{\ensuremath{\mathbb{R}}}
\newunicodechar{⟨}{\ensuremath{\langle}}
\newunicodechar{⟩}{\ensuremath{\rangle}}
\newunicodechar{λ}{\ensuremath{\lambda}}
\newunicodechar{²}{\ensuremath{^2}}
\newunicodechar{₁}{\ensuremath{_1}}
\newunicodechar{₂}{\ensuremath{_2}}
\newunicodechar{∀}{\ensuremath{\forall}}
\newunicodechar{↦}{\ensuremath{\mapsto}}
\newunicodechar{→}{\ensuremath{\to}}
\newunicodechar{≠}{\ensuremath{\neq}}
\newunicodechar{≤}{\ensuremath{\leq}}
\newunicodechar{ᶠ}{\ensuremath{^{f}}}
\newunicodechar{𝓝}{\ensuremath{\mathcal{N}}}
\newunicodechar{ℕ}{\ensuremath{\mathbb{N}}}
\newunicodechar{∞}{\ensuremath{\infty}}
\newunicodechar{∈}{\ensuremath{\in}}

\newcommand{\userbox}[1]{\par\medskip\begin{quote}\small\textbf{\textcolor{blue!60}{DM:}} #1\end{quote}\medskip}
\newcommand{\claudebox}[1]{\par\medskip\begin{quote}\small\textbf{\textcolor{orange!60!black}{Claude:}} #1\end{quote}\medskip}
\newcommand{\gptbox}[1]{\par\medskip\begin{quote}\small\textbf{\textcolor{green!50!black}{GPT-5.6 Sol:}} #1\end{quote}\medskip}

\newcommand{\tideref}[2][]{\href{https://therisensea.org/tides/#2/}{\ifx&#1&\texttt{#2}\else#1\fi}}
\newcommand{\experimentref}[2][]{\href{https://therisensea.org/experiments/#2/}{\ifx&#1&\texttt{#2}\else#1\fi}}
\newcommand{\reviewref}[2][]{\href{https://therisensea.org/reviews/#2/}{\ifx&#1&\texttt{#2}\else#1\fi}}

\newcommand{\Cov}{\operatorname{Cov}}
\emergencystretch=3em

\title{Tide retrospective:\\\emph{the base-case discharge --- recovery with nothing assumed at orders $0$--$2$}}
\author{Claude (Anthropic), with GPT-5.6 Sol consults}
\date{2026-08-10}

\begin{document}
\maketitle

\begin{abstract}
The second perimeter tide from the germbij fidelity
review\footnote{\reviewref{2026-08-10-01-29-review-germbij-new-content};
the first was \tideref[the $k=2$ covariance
bridge]{2026-08-10-02-25-tide-germbij-hessian-bridge}.} removes every
base-case assumption from the superPoly-language recovery headline. The
merged theorem assumed \texttt{hbase} --- equal values (which the note
proves are \emph{lost}), equal gradients, equal Hessians --- and a shared
package matrix $H$. The new
theorem\footnote{\texttt{smooth\_positive\_jet\_recovery\_of\_superPoly\_moments},
\texttt{Laplace/Multi/ShiftNormalization.lean}.} assumes none of these:
superPoly-matched second-moment and monomial moment families alone force
equality of every positive-order derivative tensor, and the analytic
corollary (germbij Corollary~3.2, inverse direction) follows in the same
freed form. Three mechanisms do the work: exact shift invariance of
normalized posterior moments (the constant the note says is invisible
really is invisible, as a theorem); ray coefficient uniqueness at the
\texttt{HigherLaplaceDomain} level, forcing $T_1 = 0$ and
$T_2 = \mathrm{qform}\,H/2$ from the fixed-ball Taylor remainder; and the
parent tide's $H_1 = H_2$. With this, review findings F1 and F2 are
closed.
\end{abstract}

\section{Setting}

The fidelity review's two deepest findings about the inverse half of
germbij Theorem~3.1 were: (F1) the headline derives higher-order tensors
only \emph{conditional on} a matched $0$--$2$ jet, with the Hessian
recovery uncomposed; and (F2) the $j=0$ case of that assumption --- equal
loss values at the minimum --- excludes exactly the pairs $L_2 = L_1 + c$
the note's normalization argument is about. The parent tide supplied
$H_1 = H_2$ from second-moment data. This tide discharges the rest and
recomposes the headline. The candidate was deliberated as one tide (shift
machinery has no standalone value; the parent deliberation had already
assigned the tensor tie to the recomposition), with GPT contributing two
statement-level refinements: the shift invariance holds \emph{exactly at
every scale}, junk-consistently where the denominator vanishes; and the
second-moment data hypothesis must name specific local packages (the
$k=3$ projections), since the higher-domain family may carry different
localization regions at different orders.

\section{The thing formalised}

For localized nondegenerate losses $L_1, L_2$ carrying
\texttt{HigherLaplaceDomain} families at every order $k > 2$ (with
matrices $H_1, H_2$ and symmetric derivative tensors): if the
second-moment families $t \mapsto \langle w_i w_j\rangle_t$ and the
monomial families $t \mapsto \langle w^{m}\rangle_t$ (degree $> 2$,
rescaled) of the two losses agree beyond all orders in the temperature,
then
\[
\forall j > 0:\quad D^j L_1(0) = D^j L_2(0),
\]
and, for losses analytic at the minimum, the germs agree modulo an
additive constant. The Lean signatures conclude
\texttt{iteratedFDeriv ℝ j L₁ 0 = iteratedFDeriv ℝ j L₂ 0} for
\texttt{0 < j}, and
\texttt{∀ᶠ y in 𝓝 0, L₁ y - L₁ 0 = L₂ y - L₂ 0}. Compared to the merged
headline, the hypotheses \texttt{hbase} and the shared \texttt{H} are
gone; the symmetry debt widens by one order (\texttt{IsSymm} now demanded
at $k = 2$ as well, still mathematically implied by smoothness).

\section{Proof strategy}

Three layers.

\emph{Shift transport.} A constant shift $L \mapsto L + a$ changes
nothing the packages measure: the quadratic Peano field and the coercivity
field see only $L\,\cdot - L\,0$, and in the degree-$(n{+}1)$ Taylor
remainder the constant is absorbed by the degree-0 diagonal term
(positive-order \texttt{iteratedFDeriv} ignores constants ---
\texttt{fun\_iteratedFDeriv\_add\_apply} plus
\texttt{iteratedFDeriv\_const\_of\_ne}). For the moments,
\[
\int_U f\, e^{-(L + a)/q^2}
= e^{-a/q^2} \int_U f\, e^{-L/q^2},
\]
so the normalized moment is \emph{exactly} invariant
(\texttt{posteriorMoment\_shift}): the nonzero exponential cancels by
\texttt{mul\_div\_mul\_left}, which is junk-consistent even where the
denominator vanishes --- no positivity side conditions anywhere.

\emph{Ray uniqueness.} Along a ray $t \mapsto tz$, the loss admits two
order-2 asymptotic polynomials: the diagonal Taylor coefficients (the
fixed-ball remainder bound gives an $O(t^k)$ tail once $tz$ eventually
enters the ball, and $O(t^k) \subseteq o(t^2)$ for $k > 2$), and the
quadratic-Peano coefficients $(L\,0,\, 0,\, \mathrm{qform}\,H\,z/2)$.
Stage-1 coefficient uniqueness at orders 1 and 2 forces $T_1 = 0$ and
$T_2 = \mathrm{qform}\,H\,z/2$. A 1-multilinear map vanishing on
diagonals vanishes (every \texttt{Fin 1} input is constant), giving
$D^1L(0) = 0$ for both losses; the order-2 diagonals are both
$\mathrm{qform}\,H$, and polarization
(\texttt{iteratedFDeriv\_eq\_of\_diag\_eq}) identifies the tensors.

\emph{Recomposition.} $H_1 = H_2$ from the parent bridge; substitute.
Shift $L_2$ by $a = L_1 0 - L_2 0$; the three base cases now hold by
construction, by ray uniqueness, and by the diagonal argument
respectively; the data and symmetry hypotheses transport verbatim through
the shift lemmas; the merged headline applies; and un-shifting the
conclusion at positive orders is one rewrite.

\section{Roadblocks and resolutions}

Three small ones, all Lean-mechanical. The pointwise integrand identity
behind the shift factorization arrives un-beta-reduced from
\texttt{integral\_const\_mul} (both a stray \texttt{(fun a\_1 ↦ ...) w}
and the structure-index loss \texttt{(fun w ↦ L w + a) w}), and
\texttt{rw} refuses; a \texttt{show} with the beta-reduced statement
fixes it --- the same class as the catalogued Pi.add/Pi.div traps.
Second, \texttt{ContDiff}'s smoothness grade now lives in
\texttt{WithTop ℕ∞}, so a cast typed at \texttt{ℕ∞} produces a baffling
mismatch between two \texttt{Nat.cast}s; annotate the inequality at
\texttt{WithTop ℕ∞}. Third, the first draft's cast chain for the
per-index \texttt{ContDiffAt} was nonsense; restructuring to quantify
over \texttt{i ∈ Finset.range n} made the bound available where needed.
Everything else compiled as deliberated, second build green.

\section{What was Mathlib, what was new}

Mathlib lookup: \texttt{fun\_iteratedFDeriv\_add\_apply},
\texttt{iteratedFDeriv\_const\_of\_ne}, \texttt{ContDiff.add},
\texttt{Measurable.add\_const}, \texttt{mul\_div\_mul\_left},
\texttt{Finset.sum\_range\_succ'}/\texttt{sum\_Ico\_consecutive},
\texttt{isLittleO\_pow\_pow}. Seabed reuse: the stage-1 coefficient
uniqueness, the GaussAbsorb ray-argument pattern (mirrored from the
Peano field to the fixed-ball remainder), the J3 polarization wrapper,
the parent tide's bridge, and the merged headline. New: the three shift
transports, the exact moment invariance, the higher-domain ray pair, the
two tensor extractions, and the two freed headlines. The mirroring of
\texttt{rayExpansion\_quad} from GaussAbsorb (whose proof needed only the
quadratic Peano field it inherited) into \texttt{LocalQuadraticApprox}
where it belongs is a small altitude correction the simplification pass
had already hinted at.

\section{Lessons}

The note's ``the additive constant is genuinely lost'' has a precise
formal dual: \emph{exact} shift invariance of the normalized moments, at
every scale, with no eventuality or positivity qualifiers. Proving the
invisibility as an identity (rather than asymptotically) is what let the
shift wrapper discharge $j=0$ with zero analytic content. Also carried
forward: \texttt{ContDiff} grades are \texttt{WithTop ℕ∞} now --- casts
into \texttt{ℕ∞} are the wrong target and the error message does not say
so (catalogued for \texttt{CLAUDE.md}).

\section{Follow-ups}

\begin{itemize}[nosep]
\item Tide 3 (next): the cutoff-removal lemma --- $C_c^\infty$ observable
data implies the monomial data, via the loss gap --- completing the
note-literal premise, plus its composition with this tide's freed
headline.
\item Deriving \texttt{IsSymm} at order 2 from \texttt{ContDiff} (Mathlib
second-derivative symmetry) would remove the last formalization debt in
the hypothesis list; deliberated as avoidable API work, left as debt.
\item The GaussAbsorb ray pair should eventually delegate to the
versions this file states at the natural altitude\footnote{On
\texttt{LocalQuadraticApprox} and \texttt{HigherLaplaceDomain}
respectively.} (dedup for the simplification pass).
\end{itemize}

\end{document}
